% =====================================================================
%  Capítulo 5 — Conectividade e conjunto de dados (existe a partir da AP2)
%  Versão 2026.2-AP2 (17/10/2026)
% =====================================================================
\chapter{Conectividade e conjunto de dados}
\label{cap:conectividade}
\rubricaJustificativa

Este capítulo documenta o \moduloDoCiclo: como o nó da AP1 passou a
falar com o mundo e o que ele coletou na campanha. A equipe escolheu BLE
e não WiFi. Por isso algumas peças do módulo não existem aqui
(\emph{broker}, tópicos, testamento MQTT, NTP), e cada seção diz qual é o
equivalente no BLE e o que ele garante ou deixa de garantir.

% ---------------------------------------------------------------------
\section{Tecnologia de comunicação}
\label{sec:tecnologia}
\rubricaJustificativa

\begin{table}[H]
\centering
\caption{Matriz de decisão de conectividade}
\label{tab:matriz}
\scriptsize
\setlength{\tabcolsep}{3pt}
\begin{tabular}{@{}L{2.0cm}L{3.0cm}C{1.45cm}C{1.45cm}C{1.55cm}C{1.55cm}L{2.8cm}@{}}
\toprule
\textbf{Critério} & \textbf{O que o projeto exige} & \textbf{WiFi} & \textbf{BLE}
& \textbf{LoRa\-WAN} & \textbf{NB-IoT, LTE-M} & \textbf{Observação} \\
\midrule
Alcance necessário
  & cerca de \SI{10}{\metre} até o celular, com paredes (Capítulo~\ref{cap:problema})
  & \serve & \serve & \serve & \serve
  & BLE medido: 20 a \SI{25}{\metre} com várias paredes; 45 a \SI{50}{\metre} com uma (\CT{21}) \\
Energia disponível
  & tomada USB; cerca de \SI{1114}{\milli\ampere\hour}/dia (Apêndice~\ref{ap:energia})
  & \serve & \serve & \serve & \serve
  & na tomada nada restringe; na bateria, o pico de TX do WiFi
    (\SI{240}{\milli\ampere}) é quase o dobro do BLE (\SI{130}{\milli\ampere}) \\
Dado por transmissão
  & \SI{19}{\byte} a cada \SI{300}{\second} (\SI{5472}{\byte}/dia)
  & \serve & \serve & \ressalva & \serve
  & LoRaWAN: 288 mensagens/dia cabem em \SI{30}{\second} de ar em SF7
    (\SI{20.7}{\second}), não em SF12 (\SI{521}{\second}) \\
Latência aceitável
  & minutos; a planta responde em horas
  & \serve & \serve & \serve & \serve & nenhuma restringe \\
Infraestrutura no local
  & o celular de quem cuida do jardim; a casa tem roteador WiFi
  & \ressalva & \serve & \naoserve & \naoserve
  & WiFi exige gravar a senha da rede no nó; LoRa exige \emph{gateway}; NB-IoT exige chip e plano \\
\bottomrule
\end{tabular}
\fonte{Elaborada pela equipe. Correntes de pico de \citeonline{espressif2023}; tempo no ar de LoRaWAN calculado para \SI{32}{\byte} no ar (19 de dados e 13 de cabeçalho), \SI{125}{\kilo\hertz}, CR 4/5.}
\end{table}

\begin{table}[H]
\centering
\caption{Marcas que mudaram desde a AP1}
\label{tab:mudancas}
\footnotesize
\begin{tabular}{@{}L{2.4cm}L{2.2cm}L{2.2cm}L{7.2cm}@{}}
\toprule
\textbf{Tecnologia} & \textbf{Na AP1 (19/09)} & \textbf{Agora} & \textbf{O que mudou, com o número} \\
\midrule
BLE & não avaliada & escolhida &
  O contrato BLE fechou o pacote em \SI{19}{\byte}, que cabe numa
  notificação com o MTU padrão (\DP{7}). Os ensaios com o nRF Connect
  mostraram que o central é o celular, e que a luz acumulada atravessa uma
  desconexão de \SI{55}{\minute} sem perda (\DP{8}). O alcance medido, de
  20 a \SI{25}{\metre} com paredes, é o dobro do que o jardim exige. \\
WiFi & candidata (AP1, Seção~7.2) & \ressalva &
  Na AP1 a equipe apontou WiFi porque o ESP32 já o tem. O produto, porém,
  é um aplicativo de celular: com WiFi o nó precisaria da senha da rede e
  de um servidor para chegar ao celular, peças que o BLE dispensa. \\
LoRaWAN & não avaliada & \naoserve &
  A conta do tempo no ar mostra que o volume cabe em SF7, mas o kit não tem
  módulo LoRa e não há \emph{gateway} perto; o alcance de quilômetros não
  é exigido por um jardim doméstico. \\
\bottomrule
\end{tabular}
\fonte{Elaborada pela equipe.}
\end{table}

\begin{decisao}{Tecnologia de comunicação}
\alternativas{WiFi, BLE, LoRa/LoRaWAN e NB-IoT/LTE-M.}
\criterio{Distância entre o nó e o primeiro equipamento que recebe o dado;
  orçamento de energia do Apêndice~\ref{ap:energia}; bytes e frequência das
  mensagens; infraestrutura existente no local.}
\escolha{BLE, com o nó como periférico GATT e o celular (no produto) ou o
  gravador (na campanha) como central.}
\evidencia{O jardim exige cerca de \SI{10}{\metre} até o celular, e o
  enlace alcançou 20 a \SI{25}{\metre} dentro de casa, com várias paredes,
  e 45 a \SI{50}{\metre} com uma parede só (\CT{21}). O nó entrega
  \SI{19}{\byte} a cada \SI{300}{\second}; na campanha, a completude foi de
  97,6\,\% (Seção~\ref{sec:qualidade}). O celular já está na casa
  do usuário: o BLE não exige roteador, senha de rede nem servidor, e o
  pico de transmissão é de \SI{130}{\milli\ampere}, contra
  \SI{240}{\milli\ampere} do WiFi \cite{espressif2023} (\folhadedados).}
\consequencia{O dado só sai do nó quando há um central por perto e ligado.
  Sem ele, temperatura e umidade se perdem (\DP{8}); não há como consultar
  o jardim de longe. A escolha deixa de valer se o usuário quiser
  acompanhar as plantas fora de casa, ou se o nó for para um canteiro a
  mais de cerca de \SI{20}{\metre} de onde o celular fica, atrás de
  paredes: aí o WiFi, com o roteador que a casa já tem, passa a ser a
  tecnologia certa, com o custo de provisionar a senha da rede no nó.}
\end{decisao}

% ---------------------------------------------------------------------
\section{Caminho do dado, do sensor ao arquivo}
\label{sec:caminho}
\rubricaFuncionalidade

\begin{figure}[H]
\centering
\begin{tikzpicture}[font=\sffamily\scriptsize,
  peca/.style={draw, rounded corners=2pt, align=center, minimum height=1.1cm, minimum width=2.3cm, fill=white, inner sep=3pt},
  seta/.style={-{Stealth[length=2mm]}, thick},
  rot/.style={font=\sffamily\tiny, align=center}]
  \node[peca] (sen) at (0,0)    {DHT11\\BH1750};
  \node[peca] (no)  at (3.5,0)  {nó ESP32\\no quarto};
  \node[peca] (pc)  at (8.6,0)  {adaptador Bluetooth\\do PC (BlueZ)};
  \node[peca] (gr)  at (12.6,0) {\texttt{gravador.py}};
  \node[peca] (csv) at (8.6,-2.6)  {\texttt{mensagens.csv}\\\texttt{eventos.csv}};
  \node[peca] (q)   at (3.5,-2.6)  {\texttt{qualidade.py}\\\texttt{limpar.py}};
  \draw[seta] (sen) -- node[rot, above]{1 fio e\\I\textsuperscript{2}C} (no);
  \draw[seta] ([yshift=2.5mm]no.east) -- node[rot, above]{BLE \SI{2.4}{\giga\hertz}\\notificação, 19 B} ([yshift=2.5mm]pc.west);
  \draw[seta] ([yshift=-2.5mm]pc.west) -- node[rot, below]{confirmação\\\texttt{[0x02, seq]}} ([yshift=-2.5mm]no.east);
  \draw[seta] (pc) -- node[rot, above]{D-Bus\\(\texttt{bleak})} (gr);
  \draw[seta] (gr.south) |- node[rot, pos=0.25, right]{grava e\\faz \texttt{fsync}} (csv.east);
  \draw[seta] (csv) -- node[rot, above]{lê} (q);
\end{tikzpicture}
\caption{Caminho do dado durante a campanha de coleta. A seta de volta é a confirmação, escrita depois de gravar}
\label{fig:caminho}
\fonte{Elaborada pela equipe.}
\end{figure}

\begin{table}[H]
\centering
\caption{Onde roda cada peça, e o que acontece se ela parar}
\label{tab:pecas}
\footnotesize
\begin{tabular}{@{}L{2.2cm}L{3.4cm}L{2.6cm}L{5.8cm}@{}}
\toprule
\textbf{Peça} & \textbf{Onde rodou} & \textbf{Alimentação} & \textbf{Se parar, o que se perde} \\
\midrule
Nó ESP32 & num quarto, junto a uma janela aberta, com sol direto em parte do dia, a \SI{50}{\centi\metre} do chão & carregador USB \SI{5}{\volt} na tomada &
  tudo o que ele deixaria de medir, e a luz acumulada até ali; o reinício
  aparece no id do reinício \\
Enlace BLE & entre o nó e o PC, a cerca de \SI{5}{\metre} & --- &
  se cai e volta, a assinatura faz o nó medir na hora; temperatura e
  umidade do intervalo se perdem, a luz não; a sequência mostra quantas
  leituras faltaram \\
Gravador & notebook Acer Aspire A515-57 com Linux Mint 22.3 (Cinnamon), dentro da casa & tomada &
  o mesmo que o enlace: sem \emph{broker} não há quem guarde no meio do
  caminho. Quem guarda é quem recebe \\
Arquivo & disco do mesmo computador & --- &
  cada linha vai para o disco antes da confirmação; um desligamento
  perde no máximo o pacote que estava chegando, que não foi confirmado \\
\bottomrule
\end{tabular}
\fonte{Elaborada pela equipe.}
\end{table}

A frase do Lab~03, \enquote{o \emph{broker} transporta; quem guarda é
quem assina}, no BLE fica mais curta: não há quem transporte. O gravador
é ao mesmo tempo o assinante e o único lugar onde o dado existe depois de
sair do nó. Por isso ele confirma só depois do \texttt{fsync}
(\DP{6}), e por isso, durante a campanha, o computador não pode suspender
nem ficar longe do nó.

% ---------------------------------------------------------------------
\section{Contrato de dados}
\label{sec:contrato-dados}
\rubricaJustificativa

O contrato é o pacote binário de \SI{19}{\byte} da versão 2 do contrato
BLE, \emph{little-endian}, notificado na característica de leitura. As
quatro chaves da disciplina (\texttt{id}, \texttt{boot}, \texttt{seq},
\texttt{ts}) têm equivalente, mas não todas dentro do pacote: o
\texttt{id} é o endereço BLE do nó, que o gravador anota, e o
\texttt{ts} é a hora de chegada, porque o nó não tem relógio (\DP{10}).

\begin{table}[H]
\centering
\caption{Contrato de dados da telemetria (versão 2)}
\label{tab:contrato-dados}
\footnotesize
\setlength{\tabcolsep}{3.5pt}
\begin{tabular}{@{}L{3.1cm}L{1.0cm}L{1.4cm}L{1.6cm}L{1.9cm}L{4.9cm}@{}}
\toprule
\textbf{Campo} & \textbf{Bytes} & \textbf{Tipo} & \textbf{Unidade} & \textbf{Exemplo} & \textbf{Para que serve} \\
\midrule
\multicolumn{6}{@{}l}{\emph{Fora do pacote, anotados pelo gravador}} \\
\texttt{endereco}     & --- & texto & --- & \texttt{8C:94:DF:} \texttt{8F:FA:16} & o \texttt{id} da disciplina: identifica o nó \\
\texttt{ts\_recebido} & --- & real  & s (época Unix) & 1792861964 & o \texttt{ts}: hora de chegada no gravador \\
\midrule
\multicolumn{6}{@{}l}{\emph{No pacote}} \\
\texttt{seq}          & 0     & \texttt{uint8}  & --- & 7 & conta os pacotes; o salto denuncia a perda, a repetição a duplicata; volta a 0 depois de 255 \\
\texttt{id\_reinicio} & 1--2  & \texttt{uint16} & --- & \texttt{0xA1B2} & o \texttt{boot}: sorteado a cada ligamento; separa um reinício de uma perda \\
\texttt{lux\_min\_acum} & 3--6 & \texttt{uint32} & lux$\cdot$min & 150.000 & luz desde que o nó ligou; a diferença entre dois pacotes é a luz do intervalo \\
\texttt{seg\_sem\_sucesso} & 7--10 & \texttt{uint32} & s & 1.830 & tempo desde a última confirmação que o nó recebeu \\
\texttt{temp\_c}      & 11--12 & \texttt{int16}  & \SI{0.01}{\celsius} & 2.514 & temperatura; \texttt{0x8000} = sem leitura \\
\texttt{umid\_pct}    & 13--14 & \texttt{uint16} & 0,01\,\% & 5.800 & umidade relativa; \texttt{0xFFFF} = sem leitura \\
\texttt{lux}          & 15--18 & \texttt{uint32} & lux & 20.000 & luz no instante da medida; \texttt{0xFFFFFFFF} = sem leitura \\
\bottomrule
\end{tabular}
\fonte{Elaborada pela equipe, a partir do contrato BLE do README (\enquote{Contrato BLE (versão 2)}).}
\end{table}

A unidade não está no nome da chave dentro do pacote, porque o pacote não
tem nomes: ela está no contrato e, no CSV do gravador, no nome da coluna
(\texttt{temp\_c}, \texttt{umid\_pct}). O pacote de exemplo da tabela é o
primeiro dos três do README, decodificado pelo \texttt{protocolo.py}.

\begin{lstlisting}[style=terminal, breakatwhitespace=false,
  caption={Uma mensagem real de telemetria, copiada do \texttt{mensagens.csv}},
  label={lst:payload}]
ts_recebido,hora_utc,endereco,tamanho,bytes_hex,seq,id_reinicio,lux_min_acum,seg_sem_sucesso,temp_c,umid_pct,lux
1792861964.405,2026-10-24T17:12:44.405+00:00,8C:94:DF:8F:FA:16,19,16eaebee4005002c010000cc10540b710f0000,22,60394,344302,300,43.0,29.0,3953
\end{lstlisting}

\noindent\textbf{Contrato GATT.} O nó anuncia o nome \texttt{FitoSfera} e um
serviço com duas características
(Tabela~\ref{tab:topicos}). A mudança da versão 1 para a 2 foi a
confirmação da aplicação e o campo de segundos sem sucesso: na versão 1 o
\enquote{enviado} da pilha BLE contava como entrega.

% ---------------------------------------------------------------------
\section{Contrato GATT e garantias de entrega}
\label{sec:topicos}
\rubricaJustificativa

No BLE não há tópicos nem QoS. O papel do contrato de tópicos é feito pelo
contrato GATT: quais características existem, em que direção o dado anda
e quem garante o quê.

\begin{table}[H]
\centering
\caption{Contrato GATT do nó (o equivalente do contrato de tópicos MQTT)}
\label{tab:topicos}
\scriptsize
\setlength{\tabcolsep}{3pt}
\begin{tabular}{@{}L{3.6cm}L{2.0cm}L{1.8cm}L{2.2cm}L{4.4cm}@{}}
\toprule
\textbf{UUID} & \textbf{Direção} & \textbf{Proprie\-dades} & \textbf{Quem usa} & \textbf{Por quê} \\
\midrule
\texttt{f1705f3a-0001-\ldots} (serviço) & --- & --- & --- &
  UUID próprio de 128 bits, anunciado: o gravador acha o nó pelo nome ou pelo serviço \\
\texttt{f1705f3a-0002-\ldots} (leitura) & nó $\rightarrow$ central & READ, NOTIFY & gravador assina &
  notificação, e não indicação: a garantia de entrega é a confirmação da aplicação (\DP{6}); antes do primeiro envio tem 0 bytes, para não mostrar leitura velha como atual \\
\texttt{f1705f3a-0003-\ldots} (comando) & central $\rightarrow$ nó & WRITE, WRITE sem resposta & gravador escreve &
  \texttt{01} mede agora; \texttt{02} \texttt{NN} confirma o pacote NN. O gravador escreve com resposta, para saber se a confirmação chegou \\
\bottomrule
\end{tabular}
\fonte{Elaborada pela equipe, a partir do contrato BLE do README. Os UUIDs completos terminam em \texttt{-4c7e-9a5b-3d0a6e2b9c10}.}
\end{table}

\begin{description}[style=unboxed, leftmargin=0pt, font=\normalfont\bfseries]
  \item[1. Sessão persistente ou limpa, no nó e no gravador.]\hfill\\
    Não há sessão: o nó não faz \emph{bonding} e não guarda nada sobre o
    central. Cada conexão começa do zero, e é a assinatura que dispara a
    entrega do que se acumulou. O equivalente ao \emph{client ID}
    duplicado é ter dois centrais: o nó aceita um por vez, então com o
    nRF Connect conectado o gravador não acha o nó. A regra da campanha
    foi deixar só o gravador.
  \item[2. O que o nó faz com as leituras quando o central está fora de alcance.]\hfill\\
    Não há fila (\DP{8}). O nó continua medindo a cada
    \SI{300}{\second}, a sequência continua andando e a luz continua sendo
    somada. Na volta, o primeiro pacote traz a luz do período inteiro; as
    temperaturas e umidades do período se perderam, e a sequência diz
    quantas foram. Uma fila de um dia caberia em \SI{5.5}{\kilo\byte} de
    RAM e é a melhoria mais barata para a AP3.
  \item[3. Keepalive e testamento.]\hfill\\
    O keepalive do BLE é o \emph{timeout} de supervisão da conexão,
    escolhido pelo central: se nenhum pacote do enlace chega nesse tempo,
    os dois lados consideram a conexão perdida. No gravador isso vira o
    evento \texttt{desconectado}, que é o \enquote{\texttt{offline}} do
    testamento. Medido na campanha: 3,6 e \SI{5.9}{\second} entre tirar
    a energia do nó e o evento \texttt{desconectado} (\CT{18}). Como o enlace pode continuar
    de pé com o \emph{firmware} travado, o gravador tem um segundo vigia:
    \SI{630}{\second} sem pacote derrubam a conexão. O laço do nó nunca
    fica parado mais que cerca de \SI{125}{\milli\second} (DHT11 mais
    BH1750), muito abaixo de qualquer \emph{timeout} de supervisão.
  \item[4. Segurança mínima.]\hfill\\
    Nenhuma neste ciclo, e isso é uma limitação declarada
    (Capítulo~\ref{cap:conclusao}). Qualquer central ao alcance pode
    conectar, ler as leituras e escrever na característica de comando,
    inclusive confirmar um pacote que não gravou. Não há credencial no
    repositório porque não há credencial: nada de \texttt{secrets.h}. O
    caminho para a AP3 é pareamento com \emph{bonding} e permissões de
    escrita só para o central pareado.
\end{description}

% ---------------------------------------------------------------------
\section{Custo real da comunicação}
\label{sec:custo-radio}
\rubricaJustificativa

Na AP1 o rádio estava desligado, e a linha do WiFi entrou como estimativa
(pico de cerca de \SI{310}{\milli\ampere}). Na AP2 o rádio é BLE e fica
ligado o tempo todo, anunciando ou conectado. Os tempos abaixo saem do
relógio do gravador, pelo \texttt{eventos.csv}, e são calculados pelo
\texttt{qualidade.py}. A equipe não tem medidor USB, então a corrente
vem da folha de dados (pico) e da estimativa do README (média).

\begin{table}[H]
\centering
\caption{Linha do rádio: estimada na AP1 \emph{versus} medida na AP2}
\label{tab:custo-radio}
\footnotesize
\setlength{\tabcolsep}{3.5pt}
\begin{tabular}{@{}L{3.6cm}R{1.9cm}R{2.1cm}L{3.4cm}L{2.8cm}@{}}
\toprule
\textbf{Etapa} & \textbf{Estimado (AP1)} & \textbf{Medido (AP2)} & \textbf{Como foi medido} & \textbf{Origem da corrente} \\
\midrule
Busca até achar o nó & --- & \SI{0.7}{\second} & \texttt{nao\_encontrado} e \texttt{encontrado} & --- \\
Conexão (achado $\rightarrow$ conectado) & --- & \SI{1.1}{\second} & \texttt{qualidade.py}, enlace BLE & --- \\
Conexão $\rightarrow$ primeiro pacote & --- & \SI{0.3}{\second} & idem & --- \\
Reconexão depois de queda & --- & \SI{20.1}{\second} & \texttt{desconectado} $\rightarrow$ \texttt{conectado} & --- \\
Transmissão (pico) & \SI{310}{\milli\ampere} (WiFi) & \SI{130}{\milli\ampere} (BLE) & --- & \folhadedados \\
Conectado e ocioso (média) & --- & \SI{45}{\milli\ampere} & não medido & \estimativa{} (README) \\
\midrule
\textbf{Rádio por dia} & --- & \SI{1080}{\milli\ampere\hour} & \multicolumn{2}{l}{Apêndice~\ref{ap:energia}} \\
\bottomrule
\end{tabular}
\fonte{Elaborada pela equipe. Corrente de pico de \citeonline{espressif2023}.}
\end{table}

No BLE o custo não está na transmissão, que dura frações de milissegundo
por pacote de \SI{19}{\byte}, mas em manter o rádio e a CPU acordados para
anunciar e manter a conexão. Com o nó na tomada isso não restringe; com
bateria, seria a linha que dimensiona o sistema (Apêndice~\ref{ap:energia}).

% ---------------------------------------------------------------------
\section{A campanha de coleta}
\label{sec:campanha}
\rubricaRobustez

\begin{table}[H]
\centering
\caption{Ficha da campanha de coleta}
\label{tab:campanha}
\footnotesize
\begin{tabular}{@{}L{5.0cm}L{9.0cm}@{}}
\toprule
\textbf{Item} & \textbf{Valor} \\
\midrule
Local de instalação do nó                & num quarto, junto a uma janela aberta, com sol direto em parte do dia, a \SI{50}{\centi\metre} do chão \\
Início e fim (data e hora)               & 24/10/2026 14:07 a 25/10/2026 14:13 (primeira leitura às 14:07:44, última às 14:12:32) \\
Duração total                            & 24,1\,h \\
Período de amostragem                    & \SI{300}{\second} \\
Leituras esperadas (duração $\div$ período) & 290 \\
Número de nós                            & 1 \\
Onde rodou o gravador                    & notebook Acer Aspire A515-57, Linux Mint 22.3 (Cinnamon), Python 3.12.3, na tomada \\
Como a equipe acompanhou a coleta        & saída do \texttt{gravador.py} no terminal e \texttt{qualidade.py} rodado durante os ensaios e algumas vezes por dia; diário de eventos anotado pela equipe \\
Coleta real ou simulada                  & real \\
Por que esta duração basta para a AP3    & temperatura e luz têm ciclo diário: com menos de \SI{24}{\hour}, o conjunto não mostra um dia inteiro. As 24,1\,h cobrem 1 ciclo diário completo \\
\bottomrule
\end{tabular}
\fonte{Elaborada pela equipe.}
\end{table}

As leituras esperadas saem da conta, e não do arquivo: o nó deveria
entregar uma leitura a cada \SI{300}{\second} durante toda a duração
declarada. Leituras extras, disparadas pela assinatura a cada reconexão,
não entram como esperadas; o \texttt{qualidade.py} as conta à parte.

% ---------------------------------------------------------------------
\section{Qualidade do conjunto de dados}
\label{sec:qualidade}
\rubricaRobustez

Os números saem do \texttt{qualidade.py} rodado sobre o arquivo final da
campanha, e a saída completa está no Apêndice~\ref{ap:campanha}. Como
cada reconexão reinicia o relógio do ciclo (Seção~\ref{sec:firmware}), as
leituras não caem numa grade fixa; por isso a completude é calculada por
janelas: a duração é dividida em janelas de \SI{300}{\second}, e conta-se
quantas têm pelo menos uma leitura única,
\[
  \text{completude}\,[\%] = \frac{\text{janelas de \SI{300}{\second} com leitura}}
                                 {\text{leituras esperadas}} \times 100.
\]
A razão entre leituras únicas do ciclo e leituras esperadas, que é a
fórmula da disciplina, também é impressa, e as duas ficam próximas quando
não há muitas reconexões.

\begin{table}[H]
\centering
\caption{Qualidade do conjunto de dados coletado}
\label{tab:qualidade}
\footnotesize
\begin{tabular}{@{}L{4.6cm}R{2.0cm}L{7.4cm}@{}}
\toprule
\textbf{Métrica} & \textbf{Valor} & \textbf{O que explica o valor} \\
\midrule
Leituras esperadas                    & 290 & Tabela~\ref{tab:campanha} \\
Leituras recebidas (únicas)           & 288 & 287 do ciclo e 1 extra; faltam as 7 da lacuna do \CT{16}, compensadas em parte pelas leituras de reconexão \\
\quad extras, fora do ciclo           & 1 & a assinatura reinicia o ciclo, então a leitura da reconexão quase sempre toma o lugar da do ciclo; só conta como extra quando chega colada a outra \\
Faltam (a sequência pulou)            & 12 & o nó mediu e ninguém recebeu: 7 no \CT{16} (gravador parado), 1 no \CT{15} (Bluetooth desligado) e 4 a seq 1 de reinícios, medida \SI{2}{\second} após ligar, antes da reconexão \\
Repetidas (descartadas)               & 0 & nenhum pacote chegou duas vezes \\
Reinícios do nó                       & 5 & todos provocados: 4 no \CT{18} (USB arrancado) e 1 no \CT{17} (botão EN) \\
Leituras sem hora do nó               & todas & o nó não tem relógio; a hora é a de chegada (\DP{10}) \\
Segundos sem confirmação (mediana · máximo) & 300 · \SI{2138}{\second} & o máximo é o primeiro pacote depois do \CT{16}; perto de \SI{300}{\second}, a confirmação anterior chegou \\
Campos ausentes (T · UR · luz)        & 0 · 0 · 0 & falhas de sensor depois das 3 tentativas \\
Mensagens fora do contrato            & 0 & tamanho diferente de \SI{19}{\byte} \\
\midrule
\textbf{Completude}                   & \textbf{97,6\,\%} & 283 de 290 janelas com leitura; as 7 vazias são todas da lacuna do \CT{16}. Pela fórmula da disciplina, 287 de 290 (99,0\,\%) \\
\bottomrule
\end{tabular}
\fonte{Saída do \texttt{qualidade.py} sobre \texttt{dados/mensagens.csv}.}
\end{table}

\begin{table}[H]
\centering
\caption{Lacunas do conjunto de dados}
\label{tab:lacunas}
\footnotesize
\setlength{\tabcolsep}{3.5pt}
\begin{tabular}{@{}L{2.4cm}L{2.4cm}R{1.6cm}R{1.6cm}L{3.0cm}L{2.9cm}@{}}
\toprule
\textbf{Início} & \textbf{Fim} & \textbf{Duração} & \textbf{Leituras perdidas} & \textbf{Causa} & \textbf{Tratamento} \\
\midrule
24/10 15:30:42 & 24/10 16:06:21 & 35\,min\,39\,s & 7 & \CT{16}: gravador parado de propósito por \SI{35}{\minute}; a luz do intervalo chegou no pacote seguinte & mantida como lacuna \\
\bottomrule
\end{tabular}
\fonte{Elaborada pela equipe, a partir da seção \enquote{Lacunas} do \texttt{qualidade.py} e do diário do Apêndice~\ref{ap:campanha}.}
\end{table}

O \texttt{qualidade.py} já sugere a causa de cada lacuna pelo que o
pacote seguinte traz: se o id do reinício mudou, o nó reiniciou; se a
sequência pulou, o nó mediu e o enlace ou o gravador falharam; se nenhum
dos dois, o nó deixou de medir. O arquivo de eventos completa a
explicação (gravador parado, desconexão).

% ---------------------------------------------------------------------
\section{Descrição do conjunto de dados}
\label{sec:dataset}
\rubricaRobustez

O arquivo final é o limpo, \caminho{dados/mensagens_limpo.csv}, gerado do
bruto do gravador por \caminho{dados/limpar.py}. A limpeza tira as linhas
fora do contrato e as repetidas, ordena por chegada e acrescenta colunas
calculadas; nenhum valor é interpolado. O bruto
(\caminho{dados/mensagens.csv}) e os eventos (\caminho{dados/eventos.csv})
ficam no repositório ao lado dele.

\begin{table}[H]
\centering
\caption{Descrição do conjunto de dados (\texttt{dados/mensagens\_limpo.csv})}
\label{tab:dataset}
\footnotesize
\setlength{\tabcolsep}{3.5pt}
\begin{tabular}{@{}L{3.0cm}L{1.5cm}L{1.9cm}L{2.6cm}L{5.0cm}@{}}
\toprule
\textbf{Campo} & \textbf{Tipo} & \textbf{Unidade} & \textbf{Faixa observada} & \textbf{Observação} \\
\midrule
\texttt{ts} & real & s (época Unix) & 1792861664 a 1792948352 & hora de chegada no gravador, relógio do PC \\
\texttt{hora\_local} & texto & --- & 24/10 14:07:44 a 25/10 14:12:32 & a mesma hora em Teresina (UTC$-$3) \\
\texttt{id\_reinicio} & texto (hex) & --- & 6 valores & muda a cada ligamento do nó \\
\texttt{seq} & inteiro & --- & 0 a 255 & volta a 0 depois de 255 \\
\texttt{temp\_c} & real & \si{\celsius} & 22,6 a 44,4 & vazio = sem leitura; incerteza $\pm$\SI{2}{\celsius} \\
\texttt{umid\_pct} & real & \% & 26 a 63 & vazio = sem leitura; passos de 1\,\% \\
\texttt{lux} & inteiro & lux & 0 a 61\,571 & vazio = sem leitura; incerteza $\pm$20\,\% \\
\texttt{lux\_min\_acum} & inteiro & lux$\cdot$min & 0 a 1\,245\,993 & recomeça em 0 a cada reinício \\
\texttt{luz\_nova\_lux\_min} & inteiro & lux$\cdot$min & 0 a 267\,969 & diferença do acumulado; é o que se soma para a luz do dia \\
\texttt{seg\_sem\_sucesso} & inteiro & s & 2 a 2138 & desde a última confirmação recebida pelo nó \\
\texttt{reinicio} & 0/1 & --- & --- & 1 na primeira leitura de um reinício \\
\texttt{perdidas\_antes} & inteiro & leituras & 0 a 7 & medidas que o nó fez e ninguém recebeu antes desta \\
\texttt{extra} & 0/1 & --- & --- & 1 se chegou a menos de meio período da anterior \\
\texttt{lacuna\_antes\_s} & inteiro & s & 2139 (1 valor) & preenchido quando o intervalo passou de \SI{450}{\second} \\
\bottomrule
\end{tabular}
\fonte{Elaborada pela equipe.}
\end{table}

\noindent Arquivo: \caminho{dados/mensagens_limpo.csv} \quad Tamanho:
\SI{22.6}{\kilo\byte} \quad Linhas: 288 leituras (mais o cabeçalho) \quad \emph{Commit}:
\texttt{6343e00} (tag \texttt{ap2})

\noindent\textbf{O que a equipe quer que um modelo aprenda com estes dados?}
A classificar a condição do ambiente do jardim como um todo, sem uma
planta específica. É um problema de classificação: a cada leitura (ou
janela de leituras), o modelo diz se o ambiente está adequado para um
jardim em geral ou em que condição de estresse está (calor, frio, ar
seco, ar úmido, pouca luz, sol forte). Os rótulos vêm de faixas gerais
de cultivo da literatura e do registro de eventos (Apêndice~\ref{ap:campanha});
a hora do dia e a luz acumulada (com ela, a luz integrada do dia em
\si{\mol\per\square\metre}, DLI) entram como atributos, porque a mesma
leitura pode ser normal ao meio-dia e anormal à noite. A comparação
planta a planta não é papel do modelo: fica com o aplicativo, que cruza
a condição do ambiente com os valores recomendados de cada planta,
guardados num banco de dados (\DP{3}).

\noindent\textbf{O conjunto tem exemplos suficientes disso?} Para treinar,
ainda não. O que conta é o número de dias completos, porque a luz do dia
só se mede com um dia inteiro: a campanha teve uma janela contínua de
24,1\,h com 97,6\,\% de completude, mas nenhum dia de calendário com mais
de 90\,\% das 288 janelas, porque começou e terminou no meio da tarde
(117 leituras em 24/10 e 171 em 25/10). Há exemplos de sol direto (até
\SI{61571}{lux}), de calor (até \SI{44.4}{\celsius}, com o sensor ao
sol), de noite escura e de resfriamento pelo ar-condicionado (até
\SI{22.6}{\celsius}), mas cada um aparece uma vez só, e não há frio nem ar
úmido de verdade (umidade máxima de 63\,\%). O conjunto serve para validar
o caminho do dado e os atributos; a AP3 precisa de vários dias seguidos,
de preferência com o nó também fora do quarto. As mudanças bruscas
(sol direto, sombra, janela aberta) estão anotadas no registro de eventos do
Apêndice~\ref{ap:campanha}; sem esse registro, a rotulagem de novembro
vira adivinhação.
